\section{五分钟版：这篇文章到底做了什么}

如果你只有五分钟，读这一节。

\begin{keypoint}
\textbf{一句话}：Witten 把三维 Chern--Simons 规范理论量子化，得到的量子态空间随曲面变化构成一个模函子；二维共形场论的共形块正是这个模函子的取值；纽结作为 Wilson 线插入后，配分函数就是纽结多项式，而且它对三维流形的改变（手术）给出不变量。
\end{keypoint}

论证分四步，每步的“钥匙”如下：

\begin{enumerate}[leftmargin=2em]
\item \textbf{写出一个不含度规的作用量。}
Chern--Simons 作用量 $L=\frac{k}{4\pi}\int_M\Tr(A\wedge \dd A+\tfrac23 A\wedge A\wedge A)$ 只依赖定向，不依赖度规 \src{353--354}。度规不在作用量里，所以量子化得到的量自动是拓扑不变量。这是“广义协变”的一种实现，与广义相对论“对度规积分”不同。

\item \textbf{沿曲面切开，读出相空间。}
把三维流形写成 $\Sigma\times\mathbb{R}$，作用量对时间是一阶的，给出泊松括号 $\{A_i^a(x),A_j^b(y)\}=\frac{4\pi}{k}\epsilon_{ij}\delta^{ab}\delta^2(x-y)$ \src{367}。高斯定律约束 $F=0$ 是第一类约束，辛商之后相空间是“$\Sigma$ 上平坦联络模去规范变换”的有限维空间 $\cM$ \src{367}。有限维，所以量子 Hilbert 空间有限维。

\item \textbf{几何量子化：从相空间到全纯截面。}
$\cM$ 在选一个复结构 $J$ 后是 K\"ahler 簇；预量子化线丛取行列式线丛 $\cL$，量子态空间是 $\cH_\Sigma^{(J)}=H^0(\cM_J,\cL^{\otimes k})$ \src{368}。$J$ 变时这些空间构成模空间上的射影平坦向量丛，“平坦”来自共形反常的单值性，单值性正是映射类群的作用 \src{369--371}。

\item \textbf{纽结进来，共形块对上。}
Wilson 线 $W_R(C)=\Tr_R P\exp\int_C A$ 不需要度规 \src{355}。穿刺点上的高斯定律是算符值的约束，Witten 的解法不是解它，而是把表示 $R$ 也换成量子对象：用 Borel--Weil--Bott 定理把 $R$ 实现为旗流形 $G/T$ 上辛形式 $\omega_R$ 的几何量子化 \src{371--372}。于是“带标记点的平坦联络模空间”的量子化恰好给出共形块空间 \src{373}。
\end{enumerate}

\textbf{它证明了什么、没证明什么}：
\begin{itemize}[leftmargin=1.6em]
\item 证明了（或论证到物理标准）：作用量不含度规、$k$ 的整性、$\cH_\Sigma$ 有限维、$S^3$ 上配分函数复现 Jones 多项式及其两变量推广 \src{355--356}；
\item 引用了别人的定理：Atiyah--Bott 的辛商、Borel--Weil--Bott、Segal 的共形块函子；
\item 没有证明：$\cH_\Sigma$ 与共形块空间的\textbf{函子性}（§5 只写 “It is obvious at this point” \src{397}）；半形式修正导致的能级平移 $k\to k+h^\vee$（本文略去）；Verlinde 公式的严格数学证明（本文给物理推导，数学证明 1994 年才有）。
\end{itemize}

\textbf{为什么值得读}：它示范了一种今天已成标准的做法——\emph{把拓扑不变量解释为某个量子场论的配分函数}。这件事后来发展成拓扑量子场论、Reshetikhin--Turaev 不变量、以及“用物理猜定理、用数学证定理”的整个工作方式。
